\subsection{Individual LSTM and the constant-velocity check}
\label{sec:individual_lstm}
The deterministic baseline encodes seven observed displacements with a 32-unit
LSTM and rolls out twelve increments with a separate recurrent decoder. A linear
head returns displacements; summation from the last visible or causally filled
position restores metric coordinates. It sees neither neighbors nor privileged
derivatives. The supplementary constant-velocity (CV) model repeats the last
observed displacement. Neither produces multiple futures.

Weights were trained from scratch on the locked manifests with Adam and squared
position error. Early stopping used clean validation ADE only (40-epoch ceiling,
eight-epoch patience). Future and neighbor mutations leave the input and output
unchanged. The common audit checked 90 prediction bundles. Implementation and
alignment were AI-assisted; the shared group loader and evaluator are referenced
in the owner provenance.

\paragraph{Accuracy and robustness.}
Table~\ref{tab:individual_lstm_matched} compares paired inputs and a single
prediction from each baseline. The reported numbers are equal-fold means after
averaging seeds. The dispersion across seeds describes training and, for noisy
conditions, the prescribed corruption draws; it is not an independent-window
confidence interval.

\begin{table}[ht]\centering\small
\begin{tabular}{lrrrr}\toprule Condition & LSTM ADE & LSTM FDE & CV ADE & CV FDE \\ \midrule
Clean & 0.574 & 1.191 & 0.534 & 1.148 \\
Noise 0.05 m & 0.791 & 1.526 & 0.900 & 1.748 \\
Noise 0.10 m & 1.127 & 2.059 & 1.434 & 2.663 \\
Noise 0.20 m & 1.871 & 3.271 & 2.607 & 4.714 \\
Missing 1 & 0.649 & 1.284 & 0.629 & 1.268 \\
Missing 2 & 0.735 & 1.383 & 0.732 & 1.393 \\
\bottomrule\end{tabular}
\caption{Matched single-forecast errors in meters, equal-fold means.}\label{tab:individual_lstm_matched}\end{table}


The two-point CV estimate gives the last observed displacement full influence
over every forecast step. That makes position noise especially consequential.
The deterministic LSTM can use the rest of the visible history, but cannot
respond to an unobserved change of destination or to a neighbor's motion.
We do not interpret the best of twenty Trajectron++ samples as a single
deterministic forecast. CPU runtime is device-specific and is excluded from any
cross-device speed ranking.
